% answers_short.tex —— 1978 年·全国卷 · 简短答案速查
% 由 tools/make_short_answers.py 自 answers.tex 的 \gksolhead/\gkanswer 提取，
% 并与 1978/1978_数学真题_全国_答案要点.md 逐题比对后生成，勿手改（改源 MD 后重跑）。

\begin{gknotice}
\textbf{使用说明}：本节只给最终答案便于快速核对，完整推导见解析版。答案为编写者独立推导，\textbf{非官方标准答案}。
\end{gknotice}

\gkpart{题目}
\begin{itemize}[leftmargin=2.6em]
  \item[\textbf{1.}] (1) $(x-2y+2z)(x-2y-2z)$；(2) $\dfrac{a^3}{4\pi}$；(3) $x\ge-1$（即 $[-1,+\infty)$）；(4) $\dfrac{\sqrt2}{2}$；(5) $\dfrac{4}{25}\sqrt b$
  \item[\textbf{2.}] $k>0$ 时为椭圆（$0<k<1$ 时长轴在 $x$ 轴，$k=1$ 时为圆 $x^2+y^2=4$，$k>1$ 时长轴在 $y$ 轴）；$k=0$ 时为两条平行直线 $y=\pm2$；$k<0$ 时为双曲线（焦点在 $y$ 轴）
  \item[\textbf{3.}] (1) 见解析版；(2) 见解析版
  \item[\textbf{4.}] $\dfrac{a+b}{2-a}$
  \item[\textbf{5.}] $B=60^\circ$，$(A,C)=(45^\circ,75^\circ)$ 或 $(75^\circ,45^\circ)$。当 $(A,C)=(45^\circ,75^\circ)$ 时 $a=8$、$b=4\sqrt6$、$c=4+4\sqrt3$；当 $(A,C)=(75^\circ,45^\circ)$ 时 $a=8$、$b=12\sqrt2-4\sqrt6$、$c=8\sqrt3-8$
  \item[\textbf{6.}] 见解析版
  \item[\textbf{7.}] (1) $m=-\dfrac54$；(2) 顶点均在直线 $L_1:y=x-\dfrac34$ 上；(3) 设直线为 $y=x+b$，则 $b\ge-1$ 时与各抛物线都相交（$b=-1$ 时相切），$b<-1$ 时不相交；相交时被每条抛物线截出的线段长均为 $2\sqrt2\sqrt{1+b}$，与 $m$ 无关
\end{itemize}
